% Mpu6500Module
\section{【中級】Mpu6500Module — I2Cバスパターン}

\begin{patternbox}{Mpu6500Module で学べるパターン}
\begin{itemize}
    \item 通信バスパターン（コンストラクタでバスポインタを受け取る）
    \item ヘッダーでのバスクラス前方宣言（\texttt{<Wire.h>}依存を.cppに閉じ込める）
    \item WHO\_AM\_I による init() でのデバイス識別
    \item 周期サンプリング + バースト読み取り
\end{itemize}
\end{patternbox}

\subsection{バスポインタの受け取り}

\begin{lstlisting}[style=cppstyle, caption={Mpu6500Module.h — バスの前方宣言とコンストラクタ}]
// TwoWireはポインタのみ使用するため前方宣言で十分
class TwoWire;

class Mpu6500Module : public IModule {
private:
    Mpu6500Config _config;
    TwoWire*      _wire;       // コンストラクタで受け取るバスポインタ
    ModuleTimer   _sampleTimer;
public:
    explicit Mpu6500Module(const Mpu6500Config& config, TwoWire* wire);
    // ...
};
\end{lstlisting}

\begin{pointbox}[ヘッダーでの前方宣言]
\texttt{class TwoWire;}と前方宣言することで、ヘッダーに\texttt{<Wire.h>}をインクルードする必要がない。
ポインタや参照の宣言には前方宣言だけで十分。
これにより、\texttt{Mpu6500Module.h}をインクルードするファイルに\texttt{Wire.h}の依存が波及しない。
\end{pointbox}

\begin{pointbox}[Configにバスポインタを含めない理由]
\texttt{Mpu6500Config}にはI2Cアドレスとピン情報だけを含め、\texttt{TwoWire*}は含めない。
バスポインタをConfigに含めると、\texttt{ProjectConfig.h}でバスインスタンスの\texttt{extern}宣言が
必要になり、インクルード順序の制約が生じる。
コンストラクタの別引数にすることで、Configは純粋な設定値のみを持つ。
\end{pointbox}

\subsection{main.cppでのバス管理}

\begin{lstlisting}[style=cppstyle, caption={main.cpp — バスの生成と初期化}]
// バスインスタンス（グローバルスコープで生成）
static TwoWire mpuWire = TwoWire(0);

// モジュール生成（Config + バスポインタ）
Mpu6500Module mpu6500Module(MPU6500_CONFIG, &mpuWire);

void setup() {
    // バス初期化は全モジュールのinit()より前に実行
    mpuWire.begin(MPU6500_CONFIG.sdaPin, MPU6500_CONFIG.sclPin);
    // ...
    mpu6500Module.init();  // init()ではbus.begin()を呼ばない
}
\end{lstlisting}

\begin{pointbox}[init()でbus.begin()を呼ばない]
バスの初期化は\texttt{main.cpp}の\texttt{setup()}で一元管理する。
モジュールの\texttt{init()}ではデバイス固有の初期化（WHO\_AM\_Iの確認、レジスタ設定等）のみ行う。
同じバスを複数モジュールが共有する場合、各モジュールが\texttt{begin()}すると二重初期化になる。
\end{pointbox}
